\documentclass[10pt,oneside]{book}

\usepackage{fontspec}
\setmainfont{Malgun Gothic}[
  ItalicFont={Malgun Gothic},
  ItalicFeatures={FakeSlant=0.2},
  BoldItalicFont={Malgun Gothic Bold},
  BoldItalicFeatures={FakeSlant=0.2}
]
\setsansfont{Malgun Gothic}
\XeTeXlinebreaklocale "ko"
\XeTeXlinebreakskip=0pt plus 1pt
\usepackage{amsmath,amssymb,mathrsfs}
\usepackage{longtable}
\usepackage[all]{xy}
\usepackage{tikz-cd}
\usepackage[a4paper,margin=22mm]{geometry}
\usepackage[
  colorlinks=true,
  linkcolor=blue,
  hypertexnames=false,
  unicode=true,
  pdfencoding=auto
]{hyperref}
\hypersetup{
  pdftitle={대수기하학 원론 IV, 제1분책},
  pdfauthor={알렉산더 그로텐디크와 장 디외도네},
  pdfsubject={EGA IV, 제1분책의 한국어 번역}
}

\renewcommand{\thesection}{\arabic{section}}
\renewcommand{\thesubsection}{\thesection.\arabic{subsection}}
\renewcommand{\bibname}{참고문헌}
\renewcommand{\contentsname}{목차}
\makeatletter
\renewcommand{\@pnumwidth}{2em}
\renewcommand{\@tocrmarg}{3em}
\makeatother
\newif\ifagoldpageseen
\agoldpageseenfalse
\newcommand{\oldpage}[2][]{%
  \ifagoldpageseen
    \ifhmode
      \vadjust pre{\penalty -10000}%
    \else
      \ifvmode\penalty -10000\fi
    \fi
  \else
    \global\agoldpageseentrue
  \fi
  \setcounter{footnote}{0}%
  \marginpar{\scriptsize\textbf{#1}~|~#2}\ignorespaces}
\newcommand{\agref}[2]{%
  \ifcsname r@#1\endcsname\hyperref[#1]{#2}\else#2\fi}
\makeatletter
\def\ag@sref@strip0.#1-#2\@nil#3{\agref{#3}{(#1)}}
\providecommand{\sref}[1]{\ag@sref@strip#1\@nil{#1}}
\makeatother
\providecommand{\sh}[1]{\mathcal{#1}}
\providecommand{\Spec}{\operatorname{Spec}}
\providecommand{\Proj}{\operatorname{Proj}}
\providecommand{\Hom}{\operatorname{Hom}}
\providecommand{\Hyp}{\operatorname{Hyp}}
\providecommand{\gr}{\operatorname{gr}}
\providecommand{\vphi}{\varphi}
\providecommand{\kres}{\kappa}
\providecommand{\isoto}{\xrightarrow{\sim}}
\providecommand{\supertilde}{^{\sim}}
\newenvironment{env}[1][]{\par\medskip\noindent\textbf{(#1)}\ }{\par}
\newenvironment{definition}[1][]{%
  \par\medskip\noindent\itshape 정의 (#1). ---\ }{\par}
\newenvironment{proposition}[1][]{%
  \par\medskip\noindent\itshape 명제 (#1). ---\ }{\par}
\newenvironment{theorem}[1][]{%
  \par\medskip\noindent\itshape 정리 (#1). ---\ }{\par}
\newenvironment{corollary}[1][]{%
  \par\medskip\noindent\itshape 따름정리 (#1). ---\ }{\par}
\newenvironment{remark}[1][]{%
  \par\medskip\noindent\itshape 주 (#1). ---\ }{\par}
\newenvironment{remarks}[1][]{%
  \par\medskip\noindent\itshape 주들 (#1). ---\ }{\par}
\newenvironment{scholium}[1][]{%
  \par\medskip\noindent\itshape 보충설명 (#1). ---\ }{\par}
\newenvironment{lemma}[1][]{%
  \par\medskip\noindent\itshape 보조정리 (#1). ---\ }{\par}
\newenvironment{example}[1][]{%
  \par\medskip\noindent\textit{예 (#1). ---\ }}{\par}
\newenvironment{examples}[1][]{%
  \par\medskip\noindent\textit{예들 (#1). ---\ }}{\par}
\newenvironment{notation}[1][]{%
  \par\medskip\noindent\textit{표기 (#1). ---\ }}{\par}
\newenvironment{proof}{%
  \par\medskip\noindent\textit{증명. ---\ }}{\par}
\newenvironment{namedtheorem}[2][]{%
  \par\medskip\noindent\itshape #2 (#1). ---\ }{\par}
\newenvironment{namedlemma}[2][]{%
  \par\medskip\noindent\itshape #2 (#1). ---\ }{\par}

